\documentclass[10pt,aspectratio=169,serif]{beamer}

\usepackage[fontset=none, scheme=plain]{ctex}
\IfFontExistsTF{Songti SC}{%
	\setCJKmainfont{Songti SC}[AutoFakeSlant=0.2]%
	\setCJKsansfont{Songti SC}[AutoFakeSlant=0.2]%
}{%
	\setCJKmainfont{FandolSong-Regular}[Extension=.otf,BoldFont=FandolSong-Bold,AutoFakeSlant=0.2]%
	\setCJKsansfont{FandolSong-Regular}[Extension=.otf,BoldFont=FandolSong-Bold]%
}
\usepackage{amsmath,amssymb,graphicx,tikz,booktabs,multicol,calligra}
\usepackage{pgfplots}
\pgfplotsset{compat=1.18}
\usepackage{tcolorbox}
\tcbuselibrary{skins}
\usepackage{XyygU}
\usepackage{caption}
\captionsetup[figure]{name=图,labelsep=period}

\usetikzlibrary{arrows.meta,positioning,calc}

% 公式间距
\setlength{\abovedisplayskip}{0pt plus 2pt minus 2pt}
\setlength{\belowdisplayskip}{0pt plus 2pt minus 2pt}

\title[习题1.1]{习题 1.1}
\subtitle{集合的概念}
\author{}
\institute{}
\date{2026-10-16}

\begin{document}

\begin{frame}
  \titlepage
\end{frame}

% ========== 第1题 ==========
\begin{frame}[t]{第1题：用属于或不属于填空}
\textbf{1.} 用符号“$\in$”或“$\notin$”填空：

\vspace{0.5em}
(1) 设 $A$ 为所有亚洲国家组成的集合，则
\[
\text{中国}\ \underline{\hspace{1em}}\ A,\ \ \text{美国}\ \underline{\hspace{1em}}\ A,\ \ \text{印度}\ \underline{\hspace{1em}}\ A,\ \ \text{英国}\ \underline{\hspace{1em}}\ A;
\]

(2) 若 $A = \{x \mid x^2 = x\}$，则 $-1$ \underline{\hspace{1em}} $A$；

(3) 若 $B = \{x \mid x^2 + x - 6 = 0\}$，则 $3$ \underline{\hspace{1em}} $B$；

(4) 若 $C = \{x \in \mathbf{N} \mid 1 \le x \le 10\}$，则 $8$ \underline{\hspace{1em}} $C$，$9.1$ \underline{\hspace{1em}} $C$。
\end{frame}

\begin{frame}[t]{第1题 解答}
\begin{tcolorbox}[colback=orange!10,colframe=orange,title=\textbf{解}]
\textbf{(1)} 中国、印度在亚洲；美国在北美洲；英国在欧洲。

所以：中国 $\in A$，美国 $\notin A$，印度 $\in A$，英国 $\notin A$。

\vspace{0.5em}
\textbf{(2)} 当 $x = -1$ 时，$(-1)^2 = 1 \ne -1$，不满足 $x^2 = x$，所以 $-1 \notin A$。

\vspace{0.5em}
\textbf{(3)} 当 $x = 3$ 时，$3^2 + 3 - 6 = 6 \ne 0$，不满足方程，所以 $3 \notin B$。

\vspace{0.5em}
\textbf{(4)} $8 \in \{1,2,\dots,10\}$，所以 $8 \in C$；$9.1$ 不是自然数，所以 $9.1 \notin C$。
\end{tcolorbox}
\end{frame}

% ========== 第2题 ==========
\begin{frame}[t]{第2题：用列举法表示集合}
\textbf{2.} 用列举法表示下列集合：

\vspace{0.5em}
(1) 大于 $1$ 且小于 $6$ 的整数；

\vspace{0.5em}
(2) $A = \{x \mid (x-1)(x+2) = 0\}$；

\vspace{0.5em}
(3) $B = \{x \in \mathbf{Z} \mid -3 < 2x - 1 < 3\}$。
\end{frame}

\begin{frame}[t]{第2题 解答}
\begin{tcolorbox}[colback=orange!10,colframe=orange,title=\textbf{解}]
\textbf{(1)} 大于 $1$ 且小于 $6$ 的整数为 $2, 3, 4, 5$，故
\[
\{2,\ 3,\ 4,\ 5\}.
\]

\vspace{0.5em}
\textbf{(2)} 由 $(x-1)(x+2) = 0$，得 $x = 1$ 或 $x = -2$，故
\[
A = \{-2,\ 1\}.
\]

\vspace{0.5em}
\textbf{(3)} 解不等式 $-3 < 2x - 1 < 3$：
\[
-2 < 2x < 4 \quad\Rightarrow\quad -1 < x < 2.
\]
又 $x \in \mathbf{Z}$，故 $x = 0, 1$，即
\[
B = \{0,\ 1\}.
\]
\end{tcolorbox}
\end{frame}

% ========== 第3题 ==========
\begin{frame}[t]{第3题：换一种方法表示集合}
\textbf{3.} 把下列集合用另一种方法表示出来：

\vspace{0.5em}
(1) $\{2,\ 4,\ 6,\ 8,\ 10\}$；

\vspace{0.5em}
(2) 由 $1$，$2$，$3$ 这三个数字抽出一部分或全部数字（没有重复）所组成的一切自然数；

\vspace{0.5em}
(3) $\{x \in \mathbf{N} \mid 3 < x < 7\}$；

\vspace{0.5em}
(4) 中国古代四大发明。
\end{frame}

\begin{frame}[t]{第3题 解答}
\begin{tcolorbox}[colback=orange!10,colframe=orange,title=\textbf{解}]
\textbf{(1)} 这是 $2$ 到 $10$ 的偶数，用描述法：
\[
\{x \mid x = 2k,\ k \in \mathbf{N}_+,\ k \le 5\}.
\]

\vspace{0.3em}
\textbf{(2)} 按位数分类列举（数字不重复）：
\begin{itemize}
  \item 一位数：$1, 2, 3$；
  \item 两位数：$12, 13, 21, 23, 31, 32$；
  \item 三位数：$123, 132, 213, 231, 312, 321$。
\end{itemize}
用列举法：$\{1, 2, 3, 12, 13, 21, 23, 31, 32, 123, 132, 213, 231, 312, 321\}$。

\vspace{0.3em}
\textbf{(3)} 满足 $3 < x < 7$ 的自然数为 $4, 5, 6$，故列举法：$\{4, 5, 6\}$。

\vspace{0.3em}
\textbf{(4)} 用列举法：$\{$造纸术，印刷术，指南针，火药$\}$。
\end{tcolorbox}
\end{frame}

% ========== 第4题 ==========
\begin{frame}[t]{第4题：用适当方法表示集合}
\textbf{4.} 用适当的方法表示下列集合：

\vspace{0.5em}
(1) 二次函数 $y = x^2 - 4$ 的函数值组成的集合；

\vspace{0.5em}
(2) 反比例函数 $y = \dfrac{2}{x}$ 的自变量的取值组成的集合；

\vspace{0.5em}
(3) 不等式 $3x \ge 4 - 2x$ 的解集。
\end{frame}

\begin{frame}[t]{第4题 解答}
\begin{tcolorbox}[colback=orange!10,colframe=orange,title=\textbf{解}]
\textbf{(1)} 函数 $y = x^2 - 4$，由于 $x^2 \ge 0$，所以 $y \ge -4$。

用描述法表示函数值集合：
\[
\{y \mid y = x^2 - 4,\ x \in \mathbf{R}\} = \{y \mid y \ge -4\}.
\]

\vspace{0.5em}
\textbf{(2)} 反比例函数 $y = \dfrac{2}{x}$ 中分母不能为 $0$，故 $x \ne 0$。

自变量取值集合：
\[
\{x \mid x \ne 0,\ x \in \mathbf{R}\}.
\]

\vspace{0.5em}
\textbf{(3)} 解不等式 $3x \ge 4 - 2x$：
\[
5x \ge 4 \quad\Rightarrow\quad x \ge \frac{4}{5}.
\]
解集为：
\[
\left\{x \mid x \ge \frac{4}{5}\right\}.
\]
\end{tcolorbox}
\end{frame}

% ========== 第5题 ==========
\begin{frame}[t]{第5题：集合论与康托尔}
\textbf{5.} 集合论是德国数学家康托尔（Georg Cantor，1845—1918）于 $19$ 世纪末创立的。当时，康托尔在解决涉及无限量研究的数学问题时，越过“数集”限制，提出了一般性的“集合”概念。关于集合论，希尔伯特赞誉其为“数学思想的惊人的产物，在纯粹理性的范畴中人类活动的最美的表现之一”，罗素描述其为“可能是这个时代所能夸耀的最伟大的工作”。

\vspace{0.5em}
\alert{请你查阅相关资料，用简短的报告阐述你对这些评价的认识。}
\end{frame}

\begin{frame}[t]{第5题 参考要点}
\begin{tcolorbox}[colback=orange!10,colframe=orange,title=\textbf{参考报告要点}]
\textbf{(1) 历史地位：} 集合论是现代数学的基石。康托尔第一次严格地处理了“无穷集合”，证明了无穷集之间也有大小之分（如有理数可列而实数不可列），打破了“整体大于部分”的传统观念。

\vspace{0.5em}
\textbf{(2) 思想价值：} 集合论用极其简单的概念（元素、属于）统一了几乎所有数学对象，成为数理逻辑、实变函数、泛函分析、拓扑等分支的共同语言。

\vspace{0.5em}
\textbf{(3) 争议与完善：} 早期集合论出现了罗素悖论，促使公理集合论（ZFC）诞生，使数学基础更加严密。希尔伯特、罗素的评价正是肯定了它在人类理性探索中的深度与美感。
\end{tcolorbox}
\end{frame}

\begin{frame}[c]
  \begin{center}
    {\Huge\calligra Thanks!}
  \end{center}
\end{frame}

\end{document}
